\lecture{Лекция 4}{}

\subsection{Многомерная минимизация: общие понятия}
\label{subsec:multidim-basics}

Задача многомерной минимизации:
\[
  \begin{cases}
    f(x) \longrightarrow \min, & x \in \R^n, \\
    x \in S .
  \end{cases}
\]

\subsubsection*{Два принципиальных подхода}

\begin{enumerate}[label=\arabic*., leftmargin=2em]
  \item \textbf{Аналитический}: решать уравнение $\nabla f(x) = 0$.
  \item \textbf{Итерационный}: от начального приближения $x_0$ строить
        последовательность $\{x_k\}$ с
        \[
          f(x_{k+1}) \le f(x_k) ,
        \]
        то есть \textbf{последовательность спуска}.
\end{enumerate}

\begin{idea}
Первый подход «более аналитичен», но уже для умеренно сложной $f$ система
$\nabla f(x) = 0$ решается не проще исходной задачи, да и даёт она лишь
стационарные точки, среди которых могут быть максимумы и сёдла. Поэтому весь
дальнейший курс — про второй подход.
\end{idea}

\subsubsection*{Характеристики алгоритмов}

\begin{definition}
Алгоритм \textbf{глобально сходится}, если последовательность, построенная по
этому алгоритму, сходится к оптимальному решению \textbf{независимо от выбора
начального приближения}.
\end{definition}

\begin{definition}
\textbf{Порядок алгоритма} — максимальный порядок производных, привлекаемых к
его реализации.
\end{definition}

Речь дальше пойдёт о методах $0$-го, $1$-го и $2$-го порядка.

\begin{definition}
\textbf{Скорость сходимости} алгоритма оценивается по величине порядка $p$:
\[
  \frac{\norm{x_{k+1} - x_*}}{\norm{x_k - x_*}^{\,p}} \le q ,
  \qquad q \in (0,1) ,
\]
причём требуется лишь, чтобы существовало $K$, начиная с которого неравенство
выполнено: $\forall k > K$.
\end{definition}

Принятые названия:
\[
  \begin{aligned}
    p &= 1 &&\text{— линейная скорость;} \\
    p &> 1 &&\text{— сверхлинейная скорость;} \\
    p &= 2 &&\text{— квадратичная скорость.}
  \end{aligned}
\]

\begin{idea}
Почему при $p = 1$ требуется $q < 1$, а при $p > 1$ это уже не нужно: при
$p=1$ неравенство означает $\norm{x_{k+1}-x_*} \le q\norm{x_k - x_*}$, то есть
геометрическую прогрессию, и сходимость обеспечивает именно $q<1$. При $p > 1$
ошибка возводится в степень, и достаточно лишь попасть в достаточно малую
окрестность $x_*$ — поэтому сверхлинейные и квадратичные методы сходятся, как
правило, только локально, зато очень быстро.
\end{idea}
